\section{Application to EO/IR domain shift}

For EO/IR geolocation and tracking, let $\Sigma_t$ represent an error covariance over position, attitude, timing, boresight, terrain intersection, or fused state. Normalize it to $Q_t$. Then:

\begin{itemize}
  \item a rotation of uncertainty axes caused by a platform maneuver can be mainly orientation change;
  \item growth of one eigenvalue due to blur, thermal noise, poor texture, or degraded GNSS is structural spectral change;
  \item two fault sequences can end with similar CEP50 yet have different recoverability and residual bias;
  \item a path metric can complement Fisher information, Cramer-Rao bounds, CEP50, and MMD-based sample shift.
\end{itemize}

\subsection{Relationship to Fisher-Rao and MMD}

\begin{table}[H]
\centering
\caption{Complementary domain-shift lenses}
\begin{tabularx}{\textwidth}{p{0.2\textwidth} X X}
\toprule
\textbf{Tool} & \textbf{What it measures} & \textbf{Engineering use} \\
\midrule
MMD \citep{gretton2012} & Difference between sample distributions in an RKHS & Detect that observed data changed \\
Fisher-Rao / Bures \citep{amari2016,hubner1992} & Geometric displacement inside a statistical model/state manifold & Measure model-relative distinguishability and uncertainty movement \\
Tangent/transverse split & Orientation change versus spectral-block redistribution & Classify benign reorientation versus structural degradation \\
Path-order metric & Non-commutativity of fault or adaptation sequences & Detect history effects hidden by similar endpoints \\
\bottomrule
\end{tabularx}
\end{table}

A useful EO/IR study would simulate identical final covariance traces produced by different sequences: blur then stabilization, stabilization then blur, timing drift then map correction, and reversed order. The question is not whether a quantum holonomy literally occurs, but whether an ordered geometric signature improves failure diagnosis or recovery prediction.

\section{Application to identity and credential fraud}

Identity fraud is path-dependent by nature. A final profile can appear individually plausible while its construction history is anomalous. Let $G_t$ be a temporal identity graph and let $Q_t$ summarize diffusion, reuse, or risk concentration across identity, device, credential, telco, document, and biometric nodes.

Candidate path-sensitive signals include:

\begin{itemize}
  \item order of address, device, SIM, and document changes;
  \item repeated transitions between enrollment and recovery channels;
  \item credential reuse followed by biometric replacement;
  \item graph-community changes before and after high-risk events;
  \item exclusion trajectories in which legitimate citizens are repeatedly rejected by an unsuitable process.
\end{itemize}

A combined risk score can be written
\[
R_t=w_sS_{\mathrm{sup}}+w_aS_{\mathrm{anom}}+w_gS_{\mathrm{graph}}+w_pS_{\mathrm{path}}+w_dS_{\mathrm{drift}}.
\]
The path component should be explainable as event sequences, graph motifs, and counterfactual order tests. It must not become an opaque justification for surveillance or exclusion.

\begin{engineerbox}
The useful lesson is not ``use quantum fraud detection.'' It is ``do not score only the final snapshot.'' Maintain a versioned event graph, compare plausible construction paths, identify non-commuting high-risk transitions, and preserve human review for adverse decisions.
\end{engineerbox}

\section{Application to STRAT-Q and Test Authority}

Two releases can have similar final dashboards while possessing very different evidence histories. A release that accumulated representative evidence progressively is not equivalent to one that reached the same KPI through late waivers, partial reruns, and non-representative environments.

Define a release state $Q_t$ from requirements, risks, test evidence, environment representativity, defect transitions, and decision confidence. Path-aware governance asks:

\begin{itemize}
  \item Did the evidence mature before or after critical design decisions?
  \item Were defects closed before or after environment changes?
  \item Did a waiver precede a test reduction, or did a test reduction create the waiver need?
  \item Were requirements changed after evidence was generated?
  \item Can the release state be replayed from versioned events?
\end{itemize}

This maps naturally to the planned OTEL spec-to-evidence backbone, XES/OCEL event models, XRAY evidence, and reproducible decision packages.

\subsection{C-level wording}

Avoid quantum terminology in executive governance. Use:

\begin{quote}
\textbf{Path-aware decision and evidence integrity:} release confidence considers not only current KPI values, but also the sequence, provenance, representativity, and reversibility of the evidence that produced them.
\end{quote}

\section{Quantum engineering learning path}

This topic is valuable for the user's three-to-five-year quantum roadmap because it connects linear algebra, Lie groups, differential geometry, quantum information, and open-system control.

\subsection{Prerequisite chain}

\begin{enumerate}[label=\arabic*.]
  \item Complex vector spaces, Hermitian matrices, eigenvalues, tensor products.
  \item Density matrices, partial trace, purification, completely positive maps.
  \item Lie groups and Lie algebras: $U(n)$, commutators, adjoint and coadjoint actions.
  \item Differential geometry: tangent spaces, bundles, connections, curvature, holonomy.
  \item Information geometry: Fisher metric, Bures metric, distinguishability.
  \item Open quantum dynamics: GKSL generators, decoherence, engineered reservoirs.
  \item Numerical methods: matrix exponentials, Lyapunov equations, polar decompositions.
\end{enumerate}

\subsection{Full practical lab}

The repository contains a first notebook. The next physical-computing lab should include:

\begin{enumerate}
  \item exact and numerical amplitude damping of a qubit;
  \item Bloch-ball trajectory and entropy evolution;
  \item tangent/transverse projection in the instantaneous eigenbasis;
  \item solution of \cref{eq:lyapunov};
  \item closed loops with discrete Uhlmann transport;
  \item a qutrit with controlled degeneracy splitting;
  \item sensitivity to ancilla noise and finite sampling;
  \item symbolic and numerical tests of proposed curvature identities.
\end{enumerate}

\section{Interactive tools}

Two standalone browser tools are included.

\subsection{Open Quantum Order Lab}

\texttt{tools/open\_quantum\_order\_lab.html} visualizes a qubit in the Bloch $x$-$z$ projection. Sliders control the initial state, rotation, damping probability, and trajectory resolution. It compares rotation-then-damping with damping-then-rotation and reports:

\begin{itemize}
  \item trace-distance order gap;
  \item final entropy in both orders;
  \item initial purity;
  \item final Bloch vectors;
  \item exportable JSON evidence.
\end{itemize}

\subsection{MMALS Path Memory Lab}

\texttt{tools/mmals\_path\_memory\_lab.html} uses a $2\times2$ covariance ellipse. Operation A rotates the representation. Operation B applies anisotropic scale and shear. It displays:

\begin{itemize}
  \item $A\rightarrow B$ and $B\rightarrow A$ ellipses;
  \item normalized endpoint gap;
  \item map commutator norm;
  \item tangent/transverse proxy ratio;
  \item entropy delta;
  \item a threshold-based MMALS action recommendation.
\end{itemize}

Both tools are dependency-free HTML files and can be opened offline.

